\scanpage{310}
\mspage{260}
que \inserted{chacun} des carrés $1(D),2(D),3(D)$ \deleted{\unsure{sont}}est $W$-cartésien\unsure{s} ssi les carrés correspondants $1(\overline D),2(\overline D)$\unsure{, --|,} $3(\overline D)$ l'est. Cela montre qu'$\unsure{\mathrm{OPS}|\mathrm{SPS}}$ $X_0\mathrel{\unsure{-|\to}} S$ dans $\mathrm{Fib}_W$. Considérons alors le diagramme \inserted{comm.} déduit de $D$
\diagram{\begin{tikzcd}[ampersand replacement=\&,column sep=small]
X'_0\arrow[d]\&{\widetilde X'_{\unsure{p|1}}\overset{\mathrm{def}}{=}X'_0\times_{X_0} X_1}\arrow[l]\arrow[d]\&X'_1\arrow[l,"\alpha"']\&X'_2\arrow[l]\arrow[d]\\
X_0\&X_1\arrow[l,arrows={-{Bar}}]\&\&X_2\arrow[ll]\arrow[ll,bend right=5,dash]
\end{tikzcd},}
où $\alpha\in W$ en vertu de l'hypothèse $W$-cartésienne sur le carré 1. On voit alors que le carré 2 est $W$-cartésien ssi le carré correspondant avec $X'_1$ remplacé par $\widetilde X'_1$ l'est, \unsure{quant} aux carrés 1, 3 ils ne changent pas par cette modification. Cela montre qu'$\mathrm{OPS}$ le carré 1 \uline{cartésien}. Considérons alors le diagramme
\diagram{\begin{tikzcd}[ampersand replacement=\&,column sep=small,row sep=small]
X'_0\arrow[dd,"{\mathrm{Fib}_W}"']\&X'_1\arrow[l]\arrow[dd,"{\mathrm{Fib}_W}"]\&{\widetilde X'_2\overset{\mathrm{def}}{=}X'_1\times_{X_{\unsure{1|0}}}X_2}\arrow[l,dash]\arrow[dd,"{\mathrm{Fib}_W}"]\&X'_2\arrow[l,"\beta"']\\
\arrow[r,phantom,"{\mathrm{cart}}" description]\&\arrow[r,phantom,"{\mathrm{cart}}" description]\&\&\\
X_0\&X_1\arrow[l]\&X_2\arrow[l]\&
\end{tikzcd}.}

Alors la prop. 19 et la \deleted{\unsure{dé}}définition montrent que le carré 2 est $W$-cartésien ssi $\beta\in W$, et \unsure{de même|donc} pour le carré 3. Ce qui achève la démonstration.

\uline{Corollaire 2} Soit
\diagram{\begin{tikzcd}[ampersand replacement=\&,row sep=small,column sep=small]
X_{\unsure{1|4}}\arrow[d,"{p_1}"']\&X_3\arrow[l,"{f_2}"']\arrow[d,"{f_1}"]\\
S\&X_2\arrow[l,"{p_2}"]
\end{tikzcd}}
un carré $W$-cartésien. \deleted{Si $p_1\in W$,} Alors
\[
\begin{aligned}
p_1\in W&\Longrightarrow f_1\in W\\
p_2\in W&\Longrightarrow f_2\in W
\end{aligned}
\]

\margin{plus précisément, si $p_{\unsure{i|j}}\in W$, alors $Q$ est $W$-cartésien ssi $f_{\unsure{i|i'}}\in W$}
% archive sheet number: 310
